\begin{center}
  \large\textbf{ABSTRAK}
\end{center}

\addcontentsline{toc}{chapter}{ABSTRAK}

\vspace{2ex}

\begingroup
% Menghilangkan padding
\setlength{\tabcolsep}{0pt}

\noindent
\begin{tabularx}{\textwidth}{l >{\centering}m{2em} X}
  Nama Mahasiswa    & : & \name{}         \\

  Judul Tugas Akhir & : & \tatitle{}      \\

  Pembimbing        & : & 1. \advisor{}   \\
                    &   & 2. \coadvisor{} \\
\end{tabularx}
\endgroup

% Ubah paragraf berikut dengan abstrak dari tugas akhir
Penelitain ini diajukan tentang simulasi tekanan darah pada pembuluh darah dengan media Virtual Reality, penelitian ini menerapkan konsep metode Network flow untuk mensimulasikan tekana darah dalam aliran darah tubuh yang akan divisualisasikan dengan media virtual reality. simulasi ini diharapkan dapat digunakan sebagai screening awal bagi penderita penyumbatan pembuluh darah dalam tubuh.


% Ubah kata-kata berikut dengan kata kunci dari tugas akhir
Kata Kunci: \emph{Virtual raelity}, Dinamika Fluida, Simulasi, Penyumbatan Pembuluh Darah
